\section{Summary}\label{summary}

\texttt{panelecon} is a comprehensive R toolkit for panel data
econometrics, consolidating a broad set of methods that are typically
scattered across multiple Stata ado-files or specialized software. The
package covers the full pipeline of modern panel analysis: pre-testing
for unit roots and parameter heterogeneity (Hsiao homogeneity tests,
Swamy parameter heterogeneity test); cross-sectional dependence testing
in quantile regressions; panel Granger causality tests including Fourier
Toda-Yamamoto and quantile Granger causality; group-mean and pooled
fully modified OLS (FMOLS) estimation; quantile regression slope
homogeneity tests; panel quantile unit root tests with common shocks and
structural breaks; fixed effects vector decomposition and filtered
estimators; and missing data detection and imputation tools for
unbalanced panels.

\section{Statement of Need}\label{statement-of-need}

Modern panel econometrics demands rigorous pre-testing for
cross-sectional dependence, slope heterogeneity, and non-stationarity
before any inference step. While packages such as \texttt{plm}
(Croissant and Millo 2008) handle standard linear panel models, advanced
tests like Fourier Toda-Yamamoto Granger causality, quantile panel unit
roots with structural breaks, and quantile common correlated effects are
unavailable in R. Researchers studying macroeconomic convergence,
energy-growth nexuses, and financial integration across countries have
had to rely on Stata commands with limited scriptability and
reproducibility. \texttt{panelecon} provides a reproducible,
script-friendly R alternative covering the pre-testing, estimation, and
inference stages of panel analysis within a single package.

\section{Usage}\label{usage}

\subsection{Pre-Testing: Heterogeneity and Cross-Sectional
Dependence}\label{pre-testing-heterogeneity-and-cross-sectional-dependence}

\begin{Shaded}
\begin{Highlighting}[]
\FunctionTok{library}\NormalTok{(panelecon)}

\CommentTok{\# Swamy parameter heterogeneity test and Hsiao homogeneity tests}
\NormalTok{result\_pre }\OtherTok{\textless{}{-}} \FunctionTok{xtpretest}\NormalTok{(y }\SpecialCharTok{\textasciitilde{}}\NormalTok{ x1 }\SpecialCharTok{+}\NormalTok{ x2, }\AttributeTok{data =}\NormalTok{ panel\_df,}
                        \AttributeTok{index =} \FunctionTok{c}\NormalTok{(}\StringTok{"id"}\NormalTok{, }\StringTok{"time"}\NormalTok{))}
\FunctionTok{print}\NormalTok{(result\_pre)}

\CommentTok{\# Cross{-}sectional dependence in quantile regression}
\NormalTok{result\_csd }\OtherTok{\textless{}{-}} \FunctionTok{xtcsdq}\NormalTok{(y }\SpecialCharTok{\textasciitilde{}}\NormalTok{ x1 }\SpecialCharTok{+}\NormalTok{ x2, }\AttributeTok{data =}\NormalTok{ panel\_df,}
                     \AttributeTok{index =} \FunctionTok{c}\NormalTok{(}\StringTok{"id"}\NormalTok{, }\StringTok{"time"}\NormalTok{), }\AttributeTok{tau =} \FloatTok{0.5}\NormalTok{)}
\FunctionTok{print}\NormalTok{(result\_csd)}
\end{Highlighting}
\end{Shaded}

\subsection{Panel Granger Causality}\label{panel-granger-causality}

\begin{Shaded}
\begin{Highlighting}[]
\CommentTok{\# Fourier Toda{-}Yamamoto and quantile panel Granger causality}
\NormalTok{result\_caus }\OtherTok{\textless{}{-}} \FunctionTok{xtpcaus}\NormalTok{(y }\SpecialCharTok{\textasciitilde{}}\NormalTok{ x, }\AttributeTok{data =}\NormalTok{ panel\_df,}
                       \AttributeTok{index =} \FunctionTok{c}\NormalTok{(}\StringTok{"id"}\NormalTok{, }\StringTok{"time"}\NormalTok{), }\AttributeTok{freq =} \DecValTok{1}\NormalTok{)}
\FunctionTok{summary}\NormalTok{(result\_caus)}
\end{Highlighting}
\end{Shaded}

\subsection{Panel FMOLS Estimation}\label{panel-fmols-estimation}

\begin{Shaded}
\begin{Highlighting}[]
\CommentTok{\# Group{-}mean and pooled FMOLS}
\NormalTok{result\_fmols }\OtherTok{\textless{}{-}} \FunctionTok{xtpcmg}\NormalTok{(y }\SpecialCharTok{\textasciitilde{}}\NormalTok{ x1 }\SpecialCharTok{+}\NormalTok{ x2, }\AttributeTok{data =}\NormalTok{ panel\_df,}
                       \AttributeTok{index =} \FunctionTok{c}\NormalTok{(}\StringTok{"id"}\NormalTok{, }\StringTok{"time"}\NormalTok{), }\AttributeTok{estimator =} \StringTok{"mg"}\NormalTok{)}
\FunctionTok{print}\NormalTok{(result\_fmols)}
\end{Highlighting}
\end{Shaded}

\subsection{Panel Quantile Unit Root
Tests}\label{panel-quantile-unit-root-tests}

\begin{Shaded}
\begin{Highlighting}[]
\CommentTok{\# Quantile unit root with common shocks and structural breaks}
\NormalTok{result\_qur }\OtherTok{\textless{}{-}} \FunctionTok{xtpqroot}\NormalTok{(y, }\AttributeTok{data =}\NormalTok{ panel\_df, }\AttributeTok{index =} \FunctionTok{c}\NormalTok{(}\StringTok{"id"}\NormalTok{, }\StringTok{"time"}\NormalTok{),}
                        \AttributeTok{tau =} \FunctionTok{c}\NormalTok{(}\FloatTok{0.25}\NormalTok{, }\FloatTok{0.5}\NormalTok{, }\FloatTok{0.75}\NormalTok{))}
\FunctionTok{summary}\NormalTok{(result\_qur)}
\end{Highlighting}
\end{Shaded}

\subsection{Missing Data Tools}\label{missing-data-tools}

\begin{Shaded}
\begin{Highlighting}[]
\CommentTok{\# Detect and impute missing data in panel}
\NormalTok{result\_mis }\OtherTok{\textless{}{-}} \FunctionTok{xtmispanel}\NormalTok{(panel\_df, }\AttributeTok{index =} \FunctionTok{c}\NormalTok{(}\StringTok{"id"}\NormalTok{, }\StringTok{"time"}\NormalTok{),}
                          \AttributeTok{method =} \StringTok{"linear"}\NormalTok{)}
\end{Highlighting}
\end{Shaded}

\section{Implementation}\label{implementation}

\texttt{panelecon} is written in pure R with optional dependencies on
\texttt{plm} and \texttt{zoo} for panel data structures. The
\texttt{xtpretest()} function implements the Swamy (1970) statistic
alongside Hsiao-type homogeneity tests. The \texttt{xtcsdq()} function
extends the Pesaran (2004) CD test to quantile regression residuals.
Fourier Toda-Yamamoto causality testing in \texttt{xtpcaus()} augments
the VAR with trigonometric terms to capture smooth structural shifts.
The \texttt{xtpcmg()} function implements mean-group and pooled FMOLS
following Pedroni (2000). Panel quantile unit root tests in
\texttt{xtpqroot()} follow the framework accommodating common factors
and structural breaks. The \texttt{xtqsh()} function implements quantile
slope homogeneity tests. Fixed effects vector decomposition is
implemented in line with Plümper and Troeger (2007).

\section*{References}\label{references}
\addcontentsline{toc}{section}{References}

\protect\phantomsection\label{refs}
\begin{CSLReferences}{1}{1}
\bibitem[\citeproctext]{ref-Croissant2008}
Croissant, Yves, and Giovanni Millo. 2008. {``Panel Data Econometrics in
{R}: The {plm} Package.''} \emph{Journal of Statistical Software} 27
(2): 1--43. \url{https://doi.org/10.18637/jss.v027.i02}.

\bibitem[\citeproctext]{ref-Pedroni2000}
Pedroni, Peter. 2000. {``Fully Modified {OLS} for Heterogeneous
Cointegrated Panels.''} \emph{Advances in Econometrics} 15: 93--130.
\url{https://doi.org/10.1016/S0731-9053(00)15004-2}.

\bibitem[\citeproctext]{ref-Pesaran2004}
Pesaran, M. Hashem. 2004. {``General Diagnostic Tests for Cross Section
Dependence in Panels.''} \emph{Cambridge Working Papers in Economics}
435. \url{https://doi.org/10.17863/CAM.5113}.

\bibitem[\citeproctext]{ref-Plumper2007}
Plümper, Thomas, and Vera E. Troeger. 2007. {``Efficient Estimation of
Time-Invariant and Rarely Changing Variables in Finite Sample Panel
Analyses with Unit Fixed Effects.''} \emph{Political Analysis} 15 (2):
124--39. \url{https://doi.org/10.1093/pan/mpm002}.

\bibitem[\citeproctext]{ref-Swamy1970}
Swamy, P. A. V. B. 1970. {``Efficient Inference in a Random Coefficient
Regression Model.''} \emph{Econometrica} 38 (2): 311--23.
\url{https://doi.org/10.2307/1913012}.

\end{CSLReferences}
